\chapter{Digital filtering, stability, and aliasing}
\label{ch:filter-theory}
The first operation in many LFP workflows is band-pass filtering.
Its interpretation requires a distinction between a sampled sequence,
its discrete-time filter, and the underlying analog field
\cite{oppenheim2010}.

\section{Poisson summation and aliasing}
For suitable rapidly decreasing $x(t)$, Fourier convention
$X(f)=\int x(t)e^{-2\pi i f t}\,dt$, and sampling frequency $F_s$,
Poisson summation gives the periodic spectrum of its samples:
\[
\sum_{n\in\mathbb Z}x(n/F_s)e^{-2\pi i n f/F_s}
=F_s\sum_{k\in\mathbb Z}X(f+kF_s).
\]
The observed spectrum is thus a \emph{sum of translated copies}.
The copies overlap unless the analog spectrum is sufficiently
bandlimited; an anti-alias filter must act \emph{before} digitization.
As a concrete counterexample, the sequences from frequencies $f$
and $f+F_s$ agree because
$e^{2\pi i(f+F_s)n/F_s}=e^{2\pi ifn/F_s}$.
No post-acquisition algorithm can distinguish them without
additional prior information.

\section{Convolution, the $z$-transform, and poles}
For a causal linear time-invariant discrete filter,
$y[n]=\sum_{k=0}^\infty h[k]x[n-k]$, with
\[
H(z)=\sum_{k=0}^\infty h[k]z^{-k}.
\]
Its frequency response is $H(e^{i\omega})$ when that boundary
value exists. Absolute summability $\sum_k|h[k]|<\infty$
is sufficient for bounded-input bounded-output (BIBO) stability.
For rational causal filters in minimal form, BIBO stability
requires all poles strictly inside the unit circle.

\begin{proposition}[First-order recursive filter]
For $y[n]=a y[n-1]+x[n]$ with zero initial condition,
$h[k]=a^k\mathbf 1_{k\ge0}$. The system is BIBO stable
if and only if $|a|<1$.
\end{proposition}
\begin{proof}
The absolute impulse-response sum is
$\sum_{k\ge0}|a|^k=(1-|a|)^{-1}$ when $|a|<1$.
If $|a|\ge1$, the impulse response is not absolutely summable;
there exists a bounded input producing an unbounded output.
\end{proof}
\section{FIR versus IIR; delay and zero-phase processing}
An FIR filter has a finite impulse response and is BIBO stable.
An IIR filter may have sharper transitions for a given order, but
its poles and finite-precision implementation need inspection.
For nonzero frequency response, define group delay
$\tau_g(\omega)=-d\arg H(e^{i\omega})/d\omega$.
A forward-backward filter ideally has response
$H(e^{i\omega})H(e^{-i\omega})=|H(e^{i\omega})|^2$
for real filter coefficients. This suppresses phase delay but
is \emph{noncausal}, changes the effective magnitude response,
and requires boundary extrapolation: edge transients are not removed.
An ideal brick-wall filter has an infinitely long sinc-shaped
impulse response; abrupt truncation creates Gibbs-type ringing.

\begin{example}[What the arsenal computes]
\path{bandpass_lfp.m} designs a Butterworth band-pass,
converts it to second-order sections and applies
\path{filtfilt}. The returned trace is a
zero-phase \emph{offline estimate}, not the output of a causal
online instrument. The initial and final samples need explicit
transient inspection. Filtering cannot remove volume conduction,
spike-sorting errors or physiological common inputs.
\end{example}

\section{Executable MATLAB experiment: two aliased sinusoids}
The following script constructs a 15\,Hz sinusoid and a 115\,Hz
sinusoid sampled at 100\,Hz. Their continuous paths are different,
but their sampled values coincide within floating-point tolerance.
The code also exports vector PDF and raster PNG figures when run
on a MATLAB-equipped workstation or self-hosted CI runner.
The source below is included \emph{directly} from the executable
file; this book does not keep a separate printed copy.
\lstinputlisting{../../examples/matlab/aliasing_demo.m}
% --- expanded edition (October 2026) ---
\section{Linear-phase FIR filters}
\begin{proposition}
If a real FIR filter of length $L$ has symmetric coefficients,
$h[k]=h[L-1-k]$, then $H(e^{i\omega})=e^{-i\omega(L-1)/2}A(\omega)$ with
$A$ real. Its group delay is the constant $(L-1)/2$ samples wherever
$A(\omega)\ne0$.
\end{proposition}
\begin{proof}
Pair the terms $k$ and $L-1-k$:
$h[k](e^{-i\omega k}+e^{-i\omega(L-1-k)})=2h[k]e^{-i\omega(L-1)/2}\cos(\omega((L-1)/2-k))$.
Summing gives the stated form with real $A$; the phase is linear with
slope $-(L-1)/2$, except for jumps of $\pi$ where $A$ changes sign.
\end{proof}
A constant group delay means every frequency component is delayed by
the same amount, so waveform shape is preserved up to a shift that can
be removed exactly by shifting the output by $(L-1)/2$ samples. This
gives a \emph{causal} alternative to forward--backward filtering for
offline analysis, with half the computational cost and without squaring
the magnitude response. The price is that the shifted output depends on
$(L-1)/2$ future samples, so it is not a real-time estimate.

\section{Butterworth filters and the bilinear transform}
An analog Butterworth low-pass filter of order $n$ and cutoff
$\Omega_c$ has $|H_a(i\Omega)|^2=1/(1+(\Omega/\Omega_c)^{2n})$: it is
maximally flat at $\Omega=0$ and falls by $3$\,dB at the cutoff. Its
poles lie equally spaced on the left half of the circle
$|s|=\Omega_c$. A digital filter is obtained by the \emph{bilinear
transform} $s=\frac2{T}\frac{1-z^{-1}}{1+z^{-1}}$, which maps the left
half-plane onto the unit disk, so stability is preserved. The map
compresses the infinite analog frequency axis into $(-\pi,\pi)$
according to $\Omega=\frac2T\tan(\omega/2)$; cutoffs must be
\emph{pre-warped} by this relation before design.

Two numerical facts matter in practice. First, high-order IIR filters
designed in transfer-function form $(b,a)$ have polynomial
coefficients whose roots are extremely sensitive to rounding, and
narrow low-frequency band-passes at high sampling rates may become
unstable or inaccurate; second-order sections (factoring into
biquads) avoid this. Second, a band-pass Butterworth of design order
$n$ has $2n$ poles, and forward--backward application doubles the
effective order and squares the magnitude response, so the $-3$\,dB
points of the designed filter become $-6$\,dB points of the
zero-phase result.

\section{Notch filters and line noise}
A second-order notch at normalized frequency $\omega_0$ places zeros on
the unit circle at $e^{\pm i\omega_0}$ and poles just inside at
$re^{\pm i\omega_0}$:
\[
H(z)=\frac{1-2\cos\omega_0z^{-1}+z^{-2}}{1-2r\cos\omega_0z^{-1}+r^2z^{-2}}.
\]
The closer $r$ is to one, the narrower the notch (bandwidth about
$(1-r)f_s/\pi$\,Hz) and the longer its impulse response, approximately
$1/(1-r)$ samples. A very narrow notch rings for a long time after a
transient, which can produce 50 or 60\,Hz artefacts after each large
spike or stimulation pulse. Alternatives include regressing out
line-frequency sinusoids estimated in sliding windows, or spectral
interpolation, each with its own assumptions.

\section{Causality, latency, and what a filter can know}
\begin{proposition}
A stable LTI filter with real impulse response is causal and has zero
phase at every frequency only if it is a pure gain, $h[n]=c\,\delta[n]$.
\end{proposition}
\begin{proof}
Zero phase with real coefficients means $H(e^{i\omega})$ is real, which
for a real impulse response is equivalent to $h[n]=h[-n]$. Causality
requires $h[n]=0$ for $n<0$, hence also for $n>0$.
\end{proof}
Every real-time filter therefore introduces delay or phase distortion.
Closed-loop experiments that trigger stimulation on an oscillatory phase
must estimate the phase from past samples only, typically by
autoregressive forward prediction or by causal filters with known group
delay, and should report the achieved phase accuracy measured offline
with a zero-phase filter on the same data.

\begin{example}[Group delay of a theta filter]
A symmetric FIR band-pass for 6--10\,Hz at $f_s=1000$\,Hz with $L=1001$
taps has group delay $500$\,ms, four theta cycles. A causal
implementation can only report the theta phase of half a second ago.
An IIR Butterworth of order 2 has much shorter but frequency-dependent
group delay, of the order of tens of milliseconds near the passband
centre and larger at the edges, so phase estimates are biased
differently for 6\,Hz and 10\,Hz components.
\end{example}

\section*{Exercises for this chapter}
\begin{exercise}
Show that an antisymmetric FIR filter, $h[k]=-h[L-1-k]$, also has
linear phase, with an extra constant phase of $\pi/2$.
\end{exercise}
\begin{exercise}
Find the pre-warped analog cutoff for a digital low-pass at 100\,Hz with
$f_s=1000$\,Hz.
\end{exercise}
\begin{exercise}
For a notch with $r=0.99$ at $f_s=1000$\,Hz, estimate the notch
bandwidth and the ringing duration.
\end{exercise}
\begin{exercise}
Show that the frequency response of a forward--backward filter is
$|H(e^{i\omega})|^2$.
\end{exercise}
